problem:  
Let $\mathbb{B}^3 \subset \mathbb{R}^3$ be the closed unit ball with boundary $\mathbb{S}^2$.  
For each smooth Jordan curve $\Gamma \subset \mathbb{S}^2$, suppose there exists a unique **area-minimizing** embedded disk $M_\Gamma \subset \mathbb{B}^3$ with $\partial M_\Gamma = \Gamma$.  
Let $\nu_\Gamma$ be the unit normal to $M_\Gamma$ and $\nu_{\mathbb{S}^2}$ the outer normal to the sphere along $\Gamma$.  

Define the boundary contact function
\[
\alpha(\Gamma) = \max_{p \in \Gamma}\angle(\nu_\Gamma(p),\nu_{\mathbb{S}^2}(p)).
\]

**Question:**  
Assume we have a smooth one-parameter family of embedded curves $\{\Gamma_t\}_{t\in[0,1)} \subset \mathbb{S}^2$ such that $\Gamma_t$ converges (in $C^2$) to a round equator $\Gamma_1$ as $t\to 1$.  
Suppose further that for each $t$, $\Gamma_t$ bounds a unique area-minimizing disk $M_t$ in $\mathbb{B}^3$.  

Is it possible that $\alpha(\Gamma_t)$ is constant in $t$ with  
\[
\alpha(\Gamma_t)\equiv\alpha_0\neq \tfrac{\pi}{2},
\]
that is, that the family $\{M_t\}$ remains embedded and area-minimizing, meeting $\mathbb{S}^2$ along $\Gamma_t$ at a uniform but non-right contact angle all the way to an equatorial limit?  

Or must any such family necessarily satisfy $\alpha_0 = \tfrac{\pi}{2}$ —i.e., does a rigidity principle force asymptotic orthogonality as $\Gamma_t \to \Gamma_1$?

---

Why is it a "good" problem:  
This formulation corrects the earlier ambiguities by (i) restricting to **area-minimizing disks**, ensuring uniqueness, (ii) making the question a **rigidity conjecture** rather than an existence claim, and (iii) distinguishing clearly between **Plateau problems** with fixed boundary and **capillary problems** with fixed contact angle.  

The question proposes a potentially deep geometric rigidity phenomenon: that all families of area-minimizing disks in the ball bounded by curves on the sphere flatten asymptotically to the equatorial plane with **forced right-angle contact**, irrespective of how their boundaries approach the equator.  
If true, this would uncover an *intrinsic geometric constraint* coupling minimality, convex boundary geometry, and boundary regularity; if false, it would yield a new mode of capillary-like minimal embeddings in the ball, currently unknown.  

Its apparent simplicity—“can the contact angle stay non-right as the boundary tends to an equator?”—masks a profound interplay of **boundary regularity theory**, **minimal surface curvature estimates**, and **variational rigidity**.  
Resolving it would likely require new techniques uniting the boundary behavior of the minimal surface equation, geometric measure theory, and oblique boundary conditions for quasilinear elliptic PDEs.  

Hence, the problem stands as a focused, conceptually simple, yet technically formidable bridge between **Plateau theory**, **capillarity**, and **geometric rigidity**, satisfying all attributes of a “good” research-level mathematical problem.